\let\negmedspace\undefined
\let\negthickspace\undefined
\documentclass[journal]{IEEEtran}
\usepackage[a5paper, margin=10mm, onecolumn]{geometry}
%\usepackage{lmodern} % Ensure lmodern is loaded for pdflatex
\usepackage{tfrupee} % Include tfrupee package

\setlength{\headheight}{1cm} % Set the height of the header box
\setlength{\headsep}{0mm}     % Set the distance between the header box and the top of the text

\usepackage{gvv-book}
\usepackage{gvv}
\usepackage{cite}
\usepackage{amsmath,amssymb,amsfonts,amsthm}
\usepackage{algorithmic}
\usepackage{graphicx}
\usepackage{textcomp}
\usepackage{xcolor}
\usepackage{txfonts}
\usepackage{listings}
\usepackage{enumitem}
\usepackage{mathtools}
\usepackage{gensymb}
\usepackage[breaklinks=true]{hyperref}
\usepackage{tkz-euclide}
\usepackage{listings}
% \usepackage{gvv}

\def\inputGnumericTable{}

\usepackage[latin1]{inputenc}

\usepackage{color}

\usepackage{array}

\usepackage{longtable}

\usepackage{calc}

\usepackage{multirow}

\usepackage{hhline}

\usepackage{ifthen}

\usepackage{lscape}
\usepackage{circuitikz}
\usepackage{comment}
\tikzstyle{block} = [rectangle, draw, fill=blue!20,
text width=4em, text centered, rounded corners, minimum height=3em]
\tikzstyle{sum} = [draw, fill=blue!10, circle, minimum size=1cm, node distance=1.5cm]
\tikzstyle{input} = [coordinate]
\tikzstyle{output} = [coordinate]
%--------------------------

\begin{document}
	
	\bibliographystyle{IEEEtran}
	\vspace{3cm}
	
	\title{12.55}
	\author{EE25BTECH11042 - Nipun Dasari}
	\maketitle
	
	\renewcommand{\thefigure}{\theenumi}
	\renewcommand{\thetable}{\theenumi}
	\setlength{\intextsep}{10pt}
	
	\numberwithin{equation}{enumi}
	\numberwithin{figure}{enumi}
	\renewcommand{\thetable}{\theenumi}
	
	\textbf{Question}:\\
	For a matrix \\
	\begin{align}
		\vec{M}=\begin{myvec}{\frac{4}{5}&\frac{3}{5}\\\frac{3}{5}&x}\end{myvec} \label{matrix}
	\end{align}
	The transpose of the matrix is equal to inverse of the matrix, i.e., 
	\begin{align}
		\vec{M}^\top = \vec{M}^{-1} \label{ortho}
	\end{align}
	The value of x is given by\\
	\solution
	From \eqref{ortho} it is evident that the matrix \eqref{matrix} is an orthogonal matrix.\\
	We can proceed by the fact that columns of an orthogonal vector are perpendicular
	\begin{align}
		\vec{M} = \begin{myvec}{\vec{c_1}&\vec{c_2}}\end{myvec}
	\end{align} 
	where
	\begin{align}
		\vec{c_1} = \begin{myvec}{\frac{4}{5}\\\frac{3}{5}}\end{myvec}\text{, }\vec{c_2} = \begin{myvec}{\frac{3}{5}\\x}\end{myvec}
	\end{align}
	By using \eqref{ortho}:
	\begin{align}
		\vec{c_1}^\top\vec{c_2} = 0\\
		\implies \begin{myvec}{\frac{4}{5}&\frac{3}{5}}\end{myvec}\begin{myvec}{\frac{3}{5}\\x}\end{myvec} = 0\\
		\implies \frac{12}{25} + \frac{3x}{5} = 0\\
		\implies x = -\frac{4}{5}
	\end{align}
\end{document}